\documentclass[12pt]{article}
\usepackage{amsmath,amssymb,amsfonts}
\usepackage[margin=1in]{geometry}

\begin{document}

\textbf{Chapter 6, Problem 26.} Solve the integral equation
\[
g(t) = \int_{-\infty}^{\infty} K(t-t')f(t')\,dt'
\]
by finding the Hilbert transform of $\hat{v}(\omega)$. Find a complex function $f(z) = u(x,y) + iv(x,y)$ such that $u(\alpha,0) = \hat{u}(\alpha)$ and $v(\alpha,0) = \hat{v}(\alpha)$, and verify that this complex function satisfies the two Hilbert transform relations (what are they?) that ensure that $u(\alpha)$ and $v(\alpha)$ are a Hilbert transform pair.

\bigskip
\textbf{Solution.}

Taking the Fourier transform of both sides of the convolution equation:
\[
\hat{g}(\omega) = \hat{K}(\omega)\hat{f}(\omega).
\]

Therefore:
\[
\hat{f}(\omega) = \frac{\hat{g}(\omega)}{\hat{K}(\omega)}.
\]

The solution is obtained by inverse Fourier transform:
\[
f(t) = \frac{1}{2\pi}\int_{-\infty}^{\infty}\frac{\hat{g}(\omega)}{\hat{K}(\omega)}e^{i\omega t}\,d\omega.
\]

Now, the connection to Hilbert transforms: if $\hat{K}(\omega)$ has certain analyticity properties (e.g., arises from a causal kernel $K(t) = 0$ for $t < 0$), then the real and imaginary parts of $\hat{K}(\omega) = \hat{K}_R(\omega) + i\hat{K}_I(\omega)$ form a Hilbert transform pair:
\[
\hat{K}_R(\omega) = \frac{1}{\pi}\mathcal{P}\int_{-\infty}^{\infty}\frac{\hat{K}_I(\omega')}{\omega'-\omega}\,d\omega',
\]
\[
\hat{K}_I(\omega) = -\frac{1}{\pi}\mathcal{P}\int_{-\infty}^{\infty}\frac{\hat{K}_R(\omega')}{\omega'-\omega}\,d\omega'.
\]

These are the Kramers--Kronig relations (dispersion relations). They hold because causality of $K(t)$ implies that $\hat{K}(\omega)$ is analytic in the upper half of the complex $\omega$-plane.

To verify: define $f(z) = u + iv$ analytic in the upper half-plane with boundary values $u(\alpha,0) = \hat{K}_R(\alpha)$ and $v(\alpha,0) = \hat{K}_I(\alpha)$. Since $f$ is analytic in the upper half-plane and vanishes at infinity, the Hilbert transform relations from Problem 24(c) apply, confirming $u$ and $v$ form a Hilbert transform pair.

\end{document}
